\chapter[Conclusão]{Conclusão} \label{capitulo16}

O presente trabalho teve como objetivo desenvolver um sistema web de apoio à avaliação psicopedagógica na educação infantil, buscando oferecer uma ferramenta para auxiliar os profissionais na organização das informações dos aprendentes, no gerenciamento de instrumentos avaliativos, no registro das aplicações e na elaboração de relatórios. A proposta partiu da necessidade de centralizar essas informações e facilitar sua consulta e acompanhamento ao longo do processo avaliativo.

O desenvolvimento da plataforma resultou em uma aplicação web com funcionalidades de autenticação e controle de acesso, cadastro de aprendentes, gerenciamento e versionamento de instrumentos, registro de aplicações e respostas e elaboração de relatórios. A estrutura implementada permite relacionar as informações das diferentes etapas da avaliação, mantendo os registros organizados em um ambiente único. Também foram desenvolvidos mecanismos de controle de acesso por perfil e propriedade dos dados, validação das informações e auditoria de operações, contribuindo para a segurança e a rastreabilidade dos registros.

Em relação aos objetivos específicos, foi possível estruturar a arquitetura da aplicação, modelar o banco de dados e implementar os módulos previstos para apoiar as atividades de avaliação psicopedagógica. O sistema também contempla recursos para preservar versões de instrumentos já utilizados em aplicações, permitindo manter a referência ao instrumento correspondente no momento da avaliação. Além disso, a elaboração de relatórios permanece sob responsabilidade do profissional, sem a realização de diagnósticos ou interpretações automáticas pelo sistema, respeitando o caráter profissional e contextualizado da avaliação psicopedagógica.

A verificação técnica foi realizada por meio de testes automatizados, que totalizaram 127 casos executados com sucesso, além da execução prática da plataforma em ambiente Docker, sem problemas que impedissem sua utilização. Complementarmente, foi aplicado um questionário com cinco profissionais da área, cujas respostas apresentaram média geral de 4,43 em uma escala de 1 a 5. Os resultados indicaram uma percepção predominantemente positiva, especialmente quanto ao cadastro e à consulta dos aprendentes, à facilidade de utilização e à organização das informações. As sugestões apresentadas pelas participantes também permitiram identificar oportunidades de aprimoramento, principalmente relacionadas à elaboração dos relatórios, ao armazenamento de documentos complementares e ao acompanhamento longitudinal dos aprendentes.

Apesar dos resultados alcançados, o desenvolvimento e a avaliação também evidenciaram limitações. Entre elas estão a integração ainda incompleta do histórico de aplicações e relatórios na página individual do aprendente, as possibilidades de aprimoramento do versionamento dos instrumentos e da auditoria, e a necessidade de ampliar os mecanismos de segurança e proteção dos dados. Além disso, a avaliação com profissionais contou com um número reduzido de participantes, de modo que seus resultados representam uma percepção inicial sobre a plataforma e não permitem generalizações para todo o público de psicopedagogos.

Como trabalhos futuros, destacam-se a implementação de uma linha do tempo para acompanhamento do histórico dos aprendentes, a inclusão de recursos para anexar documentos e registros complementares, o aprimoramento da estrutura dos relatórios e a disponibilização de funcionalidades para planejamento das intervenções. Também poderão ser ampliados os mecanismos de rastreabilidade, segurança e consulta das informações, bem como realizadas novas avaliações de usabilidade com um número maior de profissionais.

Conclui-se que o objetivo geral do trabalho foi alcançado, uma vez que foi desenvolvido um sistema web funcional que reúne recursos para apoiar a organização e a documentação das etapas da avaliação psicopedagógica na educação infantil. Embora ainda existam possibilidades de evolução, a plataforma apresenta uma base tecnológica e funcional para auxiliar os profissionais na centralização dos registros e na consulta das informações dos aprendentes, contribuindo para a organização do processo avaliativo e para a documentação de seu acompanhamento.